%! TeX program = lualatex
\documentclass[main]{subfiles}

\begin{document}

\title[AutoML: BO]{Practical Aspects of BO}
\subtitle{Design Size, Acquisition Optimization, Termination, Robustness}
\author[Lars Kotthoff]{Lars Kotthoff}
\date{}

\titlebackground*{images/AutoML_HPO_2_2}

%\institute{}
%\date{}

    \maketitle


\begin{frame}{Size of the Initial Design}
    \begin{itemize}
        \item Should not be too small
        \item Should not be too large
        \item Scale with the dimensionality of the search space
        \item Certain surrogate models may impose restrictions on the lower bound of the size of the initial design
        \item Rule of thumb: $4d$
    \end{itemize}
\end{frame}

%----------------------------------------------------------------------

\begin{frame}{Optimizing the Acquisition Function}
    \begin{itemize}
        \item Optimizing the acquisition function to find the next candidate point is comparably cheap
        \item Still can be a hard optimization problem: non-linear, multimodal
        \item Properly optimizing the acquisition function can be crucial for performance
    \end{itemize}
    \begin{center}
        \includegraphics[width=.6\textwidth]{images/practical.png}
    \end{center}
    \begin{footnotesize}\begin{center}
        Example EI landscape for a 2D problem.
    \end{center}\end{footnotesize}
\end{frame}

%----------------------------------------------------------------------

\begin{frame}{Optimizing the Acquisition Function}
    \begin{itemize}
        \item Choice of optimizer depends on the search space
            \begin{itemize}
                \item Numeric: L-BFGS-B with restarts, DIRECT, \ldots
                \item Mixed: EAs, local search with restarts, \ldots
                \item Random search can always be used but may not find a good solution (even with a large budget)
            \end{itemize}
        \item Sometimes, gradient-based optimizers can be used (e.g.\ when using a GP as SM)
    \end{itemize}
\end{frame}

%----------------------------------------------------------------------

\begin{frame}{When to Terminate}
    \begin{itemize}
        \item After a certain number of evaluations
            \begin{itemize}
                \item Potentially scaling the budget with the dimensionality of the search space
            \end{itemize}
        \item After a certain runtime
        \item Specify a target threshold (of the objective or acquisition function)
        \item Based on stagnation (of the objective or acquisition function)
    \end{itemize}
\end{frame}

%----------------------------------------------------------------------

\begin{frame}{Robustness}
    A good BO implementation should be \textbf{robust}:
    \begin{itemize}
        \item Handle crashing surrogate models
            \begin{itemize}
                \item Especially a GP can result in errors during training (kernel matrix not invertible, training points too close to each other)
                \item Use a fallback SM (e.g.\ Random Forest) or catch errors and propose a new candidate uniformly at random
            \end{itemize}
        \item Automatically detect input types (numerical, categorical, hierarchical) and choose an appropriate SM
        \item Have sensible defaults that work well for most scenarios
    \end{itemize}
\end{frame}

%----------------------------------------------------------------------

\begin{frame}{Choosing the Right Components}
    \begin{itemize}
        \item BO is modular: SM, acquisition function, acquisition function optimizer
        \item Different choices induce different overhead, e.g.\ GP vs.\ RF, entropy-based acquisition functions vs.\ cheap ones like EI, thorough acquisition optimization with a large budget
        \item Choosing the right components usually depends on the concrete application at hand and how expensive the evaluation of the black box itself is
    \end{itemize}
\end{frame}

%----------------------------------------------------------------------

\begin{frame}{Summary}
\begin{enumerate}
    \item Good initial design size scales with dimensionality; rule of thumb: $4d$
    \item Acquisition optimization is cheap but non-trivial; choose an optimizer suited to the search space
    \item Termination can be based on evaluation budget, runtime, targets, or stagnation
    \item Robust BO requires fallbacks, input-type awareness, and sensible defaults; BO is modular and component choices depend on the application
\end{enumerate}
\end{frame}

\end{document}
